\chapter{Rozkład Jordana--Chevalleya i funkcje macierzy}
\label{ch:lin18-jordan-chevalley}
\index{rozkład Jordana--Chevalleya@Rozkład Jordana--Chevalleya}
\index{rzut spektralny@Rzut spektralny}
\index{funkcja macierzy@Funkcja macierzy}
\index{wykładnicza macierzy@Wykładnicza macierzy}

\noindent\textbf{Cel rozdziału.}
W Rozdziale~\ref{ch:lin17-jordan}
uzyskaliśmy postać Jordana operatorów, których
wielomian charakterystyczny rozszczepia się nad
rozpatrywanym ciałem.
Pokażemy teraz, jak rozdzielić taki operator
na przemienne części: diagonalizowalną i nilpotentną.
Następnie wyprowadzimy wzory na wielomiany,
potęgi oraz odwrotności operatora.
Nad \(\mathbb C\) zastosujemy te konstrukcje
do funkcji analitycznych i macierzy wykładniczej.

Do końca sekcji dotyczącej potęg i odwrotności
\(K\) oznacza dowolne ciało,
\(V\) skończenie wymiarową przestrzeń wektorową
nad \(K\), a \(T\in\operatorname{End}_K(V)\).
Zakładamy, że wielomian minimalny \(m_T\)
rozszczepia się nad \(K\).
Przez \(\sigma_K(T)\) oznaczamy zbiór wartości
własnych \(T\) należących do \(K\).
Wymiar \(0\) dopuszczamy z konwencją
\(I_{\{0\}}=0\), pustymi sumami równymi \(0\)
oraz pustymi iloczynami równymi \(1\).

\section{Rzuty na uogólnione przestrzenie własne}
\label{sec:lin18-generalized-projectors}

\begin{twierdzenie}[Wielomianowe rzuty na uogólnione przestrzenie własne]
\label{thm:lin18-projectors}
Niech \(T\in\operatorname{End}_K(V)\)
ma wielomian minimalny
\[
 m_T(t)=\prod_{\lambda\in\sigma_K(T)}
             (t-\lambda)^{a_\lambda},
 \qquad a_\lambda\ge1.
\]
Dla każdego \(\lambda\in\sigma_K(T)\)
istnieje wielomian \(q_\lambda\in K[t]\),
który spełnia
\[
 \begin{split}
 q_\lambda(t)&\equiv1
        \pmod{(t-\lambda)^{a_\lambda}},\\
 q_\lambda(t)&\equiv0
        \pmod{(t-\mu)^{a_\mu}}
        \qquad(\mu\ne\lambda).
 \end{split}
\]
Niech \(P_\lambda=q_\lambda(T)\).
Operator \(P_\lambda\) jest niezależny
od wyboru takiego wielomianu.
Ponadto
\[
 \begin{gathered}
 P_\lambda^2=P_\lambda,\qquad
 P_\lambda P_\mu=0\quad(\lambda\ne\mu),\\
 \sum_{\lambda\in\sigma_K(T)}P_\lambda=I_V,\qquad
 \operatorname{im}P_\lambda
       =E_\lambda^{\mathrm{gen}}(T).
 \end{gathered}
\]
Wszystkie \(P_\lambda\) przemieniają się z \(T\).
Przy \(V=\{0\}\) rodzina rzutów jest pusta,
a jej suma to \(I_V=0\).
\end{twierdzenie}
\begin{proof}
Przy \(V=\{0\}\) wniosek jest zgodny
z podanymi konwencjami.
Załóżmy \(V\ne\{0\}\), ustalmy
\(\lambda\in\sigma_K(T)\) i oznaczmy
\[
 f_\lambda(t)=(t-\lambda)^{a_\lambda},
 \qquad
 g_\lambda(t)=m_T(t)/f_\lambda(t).
\]
Wielomiany \(f_\lambda\) i \(g_\lambda\)
są względnie pierwsze, ponieważ czynniki
liniowe odpowiadają parami różnym wartościom.
Z tożsamości Bézouta istnieją
\(u_\lambda,v_\lambda\in K[t]\)
takie, że
\[
 u_\lambda g_\lambda+v_\lambda f_\lambda=1.
\]
Wybierzmy \(q_\lambda=u_\lambda g_\lambda\).
Z tej tożsamości wynika pierwsza kongruencja.
Dla \(\mu\ne\lambda\) wielomian
\(g_\lambda\) jest podzielny
przez \(f_\mu\), co daje drugą kongruencję.

Na mocy rozkładu z
Wniosku~\ref{cor:lin16-generalized}
mamy
\[
 V=\bigoplus_{\mu\in\sigma_K(T)}W_\mu,
 \qquad
 W_\mu=E_\mu^{\mathrm{gen}}(T).
\]
Ze
Stwierdzenia~\ref{prop:lin16-gen-stabilize}
wynika \(f_\mu(T)|_{W_\mu}=0\).
Jeśli \(\mu=\lambda\), to
\(q_\lambda-1\) jest wielokrotnością
\(f_\lambda\), więc
\(q_\lambda(T)|_{W_\lambda}=I_{W_\lambda}\).
Jeśli \(\mu\ne\lambda\), to
\(q_\lambda\) jest wielokrotnością \(f_\mu\),
więc \(q_\lambda(T)|_{W_\mu}=0\).
Tak więc \(P_\lambda\) jest rzutem
na \(W_\lambda\) wzdłuż sumy wszystkich
pozostałych \(W_\mu\).
Z tego opisu bezpośrednio wynikają
idempotentność, zerowe iloczyny różnych
rzutów, ich suma równa \(I_V\)
oraz równość obrazu z \(W_\lambda\).

Dowolny drugi wielomian spełniający
te same kongruencje działa identycznie
na każdym \(W_\mu\), a więc na całym \(V\).
Wynik nie zależy od wyboru \(q_\lambda\).
Wreszcie każdy \(P_\lambda\) jest wielomianem
w \(T\), dlatego \(TP_\lambda=P_\lambda T\).
\end{proof}

\begin{uwaga}[Rzuty nie muszą być ortogonalne]
\label{rem:lin18-not-orthogonal}
Przy danym rozkładzie
\(V=\bigoplus_\lambda W_\lambda\)
rzuty \(P_\lambda\) są rzutami algebraicznymi.
Nie zakładamy istnienia iloczynu skalarnego,
a nawet w przestrzeni unitarnej nie muszą
być samosprzężone ani ortogonalne.
Dla operatorów normalnych mamy
silniejszy rozkład na rzuty ortogonalne
z Twierdzenia~\ref{thm:lin14-spectral-projections}.
\end{uwaga}

\begin{przyklad}[Rzut otrzymany z wielomianu]
\label{prz:lin18-projector-polynomial}
Niech \(K=\mathbb Q\) i rozważmy
\[
 T=
 \begin{pmatrix}
  1&1&0\\0&1&0\\0&0&3
 \end{pmatrix}.
\]
Mamy \(m_T(t)=(t-1)^2(t-3)\).
Wielomian
\[
 q_3(t)=\frac{(t-1)^2}{4}
\]
jest podzielny przez \((t-1)^2\),
a \(q_3(3)=1\), czyli
\(q_3(t)\equiv1\pmod{t-3}\).
Zatem
\[
 P_3=q_3(T)=\operatorname{diag}(0,0,1),
 \qquad
 P_1=I-P_3=\operatorname{diag}(1,1,0).
\]
Nie wyznaczaliśmy bazy wektorów własnych:
rzuty uzyskaliśmy bezpośrednio
z wielomianów w macierzy \(T\).
\end{przyklad}

\section{Rozkład Jordana--Chevalleya}
\label{sec:lin18-chevalley}

\begin{definicja}[Część diagonalizowalna i nilpotentna]
\label{def:lin18-chevalley}
Jeśli \(T=S+N\), gdzie
\(S,N\in\operatorname{End}_K(V)\),
\(S\) jest diagonalizowalny nad \(K\),
\(N\) jest nilpotentny oraz \(SN=NS\),
to taką parę nazywamy
\emph{rozkładem Jordana--Chevalleya operatora \(T\)}
nad \(K\).
\end{definicja}

\begin{twierdzenie}[Istnienie i jednoznaczność rozkładu Jordana--Chevalleya]
\label{thm:lin18-chevalley}
Niech \(K\) będzie dowolnym ciałem,
\(V\) przestrzenią skończenie wymiarową nad \(K\),
a \(T\in\operatorname{End}_K(V)\).
Jeżeli \(m_T\) rozszczepia się
nad \(K\) na czynniki liniowe, to istnieje
dokładnie jedna para operatorów \(S,N\)
spełniająca
\[
 T=S+N,\qquad SN=NS,
\]
przy czym \(S\) jest diagonalizowalny
nad \(K\), a \(N\) jest nilpotentny.
Oba operatory są wielomianami w \(T\).
Dokładniej, dla rzutów
z Twierdzenia~\ref{thm:lin18-projectors}
\[
 S=\sum_{\lambda\in\sigma_K(T)}
       \lambda P_\lambda,\qquad N=T-S.
\]
Ponadto \(T\) jest diagonalizowalny
dokładnie wtedy, gdy \(N=0\).
\end{twierdzenie}
\begin{proof}
Z rozkładu na uogólnione przestrzenie własne
wynika
\[
 V=\bigoplus_{\lambda}W_\lambda,
 \qquad W_\lambda=\operatorname{im}P_\lambda.
\]
Zdefiniujmy \(S=\sum_\lambda\lambda P_\lambda\)
oraz \(N=T-S\).
Każdy składnik \(W_\lambda\)
jest niezmienniczy względem \(S\) i \(N\);
operator \(S\) działa na nim
jako \(\lambda I\).
Łącząc dowolne bazy przestrzeni
\(W_\lambda\), otrzymujemy bazę,
w której \(S\) jest diagonalny.
Ponieważ wszystkie \(P_\lambda\)
są wielomianami w \(T\),
również \(S\) i \(N\) są wielomianami
w \(T\), a zatem \(SN=NS\).

Na \(W_\lambda\) operator \(N\)
jest równy \(T-\lambda I\).
Z definicji uogólnionej przestrzeni własnej
i stabilizacji jąder istnieje
liczba \(a_\lambda\ge1\) taka, że
\[
 N^{a_\lambda}|_{W_\lambda}=0.
\]
Jeśli \(V\ne0\), wybierając
\(a=\max_\lambda a_\lambda\),
otrzymujemy \(N^a=0\) na każdym składniku,
a więc na całym \(V\).
Gdy \(V=0\), mamy \(N=0\), zatem
nilpotentność jest również spełniona.

Udowodnimy jednoznaczność.
Załóżmy, że \(T=S'+N'\), gdzie
\(S'\) jest diagonalizowalny nad \(K\),
\(N'\) jest nilpotentny oraz \(S'N'=N'S'\).
Niech \(E_\mu(S')\) będzie przestrzenią
własną \(S'\) dla wartości \(\mu\in K\).
Komutowanie zapewnia, że
\(N'(E_\mu(S'))\subset E_\mu(S')\).
Na tej przestrzeni
\[
 (T-\mu I)|_{E_\mu(S')}=N'|_{E_\mu(S')}.
\]
Ponieważ \(N'\) jest nilpotentny,
każdy wektor z \(E_\mu(S')\)
należy do \(E_\mu^{\mathrm{gen}}(T)\).
Ponadto
\[
 V=\bigoplus_{\mu\in\sigma_K(S')}E_\mu(S'),
\]
bo \(S'\) jest diagonalizowalny.
Jeżeli \(\mu\) nie jest wartością własną \(T\),
to \(E_\mu^{\mathrm{gen}}(T)=\{0\}\),
a zatem żaden niezerowy wektor
nie może należeć do \(E_\mu(S')\).
Z kolei dla \(\lambda\in\sigma_K(T)\)
każdy \(w\in W_\lambda\) rozkłada się
na wektory ze wspomnianych
przestrzeni własnych \(S'\).
Każdy taki składnik należy również
do pewnej przestrzeni \(W_\mu\).
Bezpośredniość sumy \(\bigoplus_\mu W_\mu\)
wymusza zerowanie wszystkich składników
dla \(\mu\ne\lambda\).
Stąd \(W_\lambda=E_\lambda(S')\).
Operator \(S'\) działa na \(W_\lambda\)
jako \(\lambda I\), czyli tak samo jak \(S\).
Ponieważ suma przestrzeni \(W_\lambda\)
jest całym \(V\), mamy \(S'=S\);
następnie \(N'=T-S'=N\).

Jeżeli \(N=0\), to \(T=S\)
jest diagonalizowalny.
Odwrotnie, jeżeli \(T\) jest diagonalizowalny,
para \((T,0)\) spełnia warunki rozkładu.
Z jego jednoznaczności wynika \(S=T\) i \(N=0\).
\end{proof}

\begin{przyklad}[Część diagonalizowalna i nilpotentna]
\label{prz:lin18-chevalley-example}
Dla operatora z
Przykładu~\ref{prz:lin18-projector-polynomial}
otrzymujemy
\[
 S=P_1+3P_3
   =\operatorname{diag}(1,1,3),
 \qquad
 N=T-S=
 \begin{pmatrix}
 0&1&0\\0&0&0\\0&0&0
 \end{pmatrix}.
\]
Mamy \(N^2=0\), \(SN=NS\),
a \(S\) jest diagonalny.
Rozkład jest więc rozkładem
Jordana--Chevalleya.
\end{przyklad}

\section{Wielomiany operatora, jego potęgi i odwrotność}
\label{sec:lin18-polynomial-calculus}

\begin{definicja}[Współczynniki rozwinięcia wielomianu]
\label{def:lin18-taylor-coefficients}
Niech \(p\in K[t]\), \(\lambda\in K\).
Dla \(j\ge0\) oznaczmy przez
\(p^{[j]}(\lambda)\in K\)
współczynnik przy \(u^j\)
w wielomianie \(p(\lambda+u)\):
\[
 p(\lambda+u)
   =\sum_{j=0}^{\deg p}p^{[j]}(\lambda)u^j.
\]
Przy \(p=0\) wszystkie te współczynniki
są równe zero.
Gdy \(K\) ma charakterystykę \(0\),
zachodzi \(p^{[j]}(\lambda)=p^{(j)}(\lambda)/j!\).
Definicja współczynnikowa
jest jednak poprawna nad każdym ciałem,
również w dodatniej charakterystyce.
\end{definicja}

\begin{twierdzenie}[Rachunek wielomianowy na uogólnionych przestrzeniach własnych]
\label{thm:lin18-polynomial}
Niech \(T=S+N\) będzie rozkładem
z Twierdzenia~\ref{thm:lin18-chevalley},
a \(P_\lambda\) odpowiadającymi mu rzutami.
Niech \(a_\lambda\) będzie wykładnikiem
czynnika \((t-\lambda)\) w \(m_T\).
Dla dowolnego wielomianu \(p\in K[t]\)
zachodzi
\[
 p(T)=
 \sum_{\lambda\in\sigma_K(T)}
 \sum_{j=0}^{a_\lambda-1}
 p^{[j]}(\lambda)\,N^jP_\lambda.
\]
W szczególności \(p(T)\) zależy
jedynie od współczynników rozwinięć
\(p(\lambda+u)\) do stopnia
mniejszego niż \(a_\lambda\)
dla poszczególnych wartości własnych.
\end{twierdzenie}
\begin{proof}
Niech \(W_\lambda=\operatorname{im}P_\lambda\).
Na tej przestrzeni \(T=\lambda I+N\),
przy czym \(N^{a_\lambda}|_{W_\lambda}=0\).
Ponieważ w pierścieniu przemiennym \(K[u]\)
mamy
\(p(\lambda+u)=\sum_jp^{[j]}(\lambda)u^j\),
po zastąpieniu \(u\) przez
nilpotentny operator \(N|_{W_\lambda}\)
otrzymujemy
\[
 p(T)|_{W_\lambda}
 =\sum_{j=0}^{a_\lambda-1}
       p^{[j]}(\lambda)\,N^j|_{W_\lambda}.
\]
Dla każdego \(v\in V\)
rozkład na rzuty daje
\(v=\sum_\lambda P_\lambda v\).
Zastosowanie \(p(T)\) do tego rozkładu
i powyższej równości na poszczególnych
\(W_\lambda\) dowodzi wzoru na \(p(T)\).
Przy \(V=0\) obie strony są zerowym
endomorfizmem.
\end{proof}

\begin{wniosek}[Potęgi i odwrotność]
\label{cor:lin18-powers-inverse}
W założeniach
Twierdzenia~\ref{thm:lin18-polynomial}
dla każdego \(k\ge0\) zachodzi
\[
 T^k=
 \sum_{\lambda\in\sigma_K(T)}
 \sum_{j=0}^{\min(k,a_\lambda-1)}
 \binom{k}{j}\lambda^{\,k-j}N^jP_\lambda,
\]
gdzie współczynniki dwumianowe
są interpretowane w ciele \(K\).
Przy \(k=j=0\) używamy zwykłej
konwencji \(0^0=1\) dla wartości
wielomianu stałego \(t^0=1\).

Jeśli \(0\notin\sigma_K(T)\),
to \(T\) jest odwracalny i
\[
 T^{-1}=
 \sum_{\lambda\in\sigma_K(T)}
 \sum_{j=0}^{a_\lambda-1}
 (-1)^j\lambda^{-j-1}N^jP_\lambda.
\]
W wymiarze \(0\) pusta suma
jest jedynym endomorfizmem,
który jest również \(I_V=T^{-1}\).
\end{wniosek}
\begin{proof}
Stosujemy
Twierdzenie~\ref{thm:lin18-polynomial}
do \(p(t)=t^k\).
Ze wzoru dwumianowego
\[
 (\lambda+u)^k
 =\sum_{j=0}^k\binom{k}{j}\lambda^{k-j}u^j
\]
otrzymujemy pierwszą tożsamość,
a wyrazy z \(j\ge a_\lambda\)
zerują się na \(W_\lambda\).

Jeżeli \(\lambda\ne0\),
to na \(W_\lambda\)
operator \(T=\lambda I+N\)
ma odwrotność
\[
 (\lambda I+N)^{-1}
 =\lambda^{-1}
   \sum_{j=0}^{a_\lambda-1}
       (-\lambda^{-1}N)^j.
\]
Mnożąc przez \(\lambda I+N\),
otrzymujemy szereg teleskopowy
\[
 (\lambda I+N)\lambda^{-1}
 \sum_{j=0}^{a_\lambda-1}(-\lambda^{-1}N)^j
 =I-(-\lambda^{-1}N)^{a_\lambda}=I,
\]
ponieważ \(N^{a_\lambda}=0\)
na tej podprzestrzeni.
Analogiczne mnożenie z drugiej strony
daje identyczność; czynniki przemieniają się.
Suma tych odwrotności na wszystkich
niezmienniczych składnikach \(W_\lambda\)
daje wzór globalny.
\end{proof}

\begin{przyklad}[Potęgi bloku Jordana]
\label{prz:lin18-jordan-powers}
Dla dowolnego ciała \(K\)
i \(J=J_3(1)\) mamy \(J=I+N\), gdzie
\[
 N=\begin{pmatrix}0&1&0\\0&0&1\\0&0&0\end{pmatrix},
 \qquad N^3=0.
\]
Zatem dla każdego \(k\ge0\)
\[
 J^k=I+kN+\binom{k}{2}N^2
 =\begin{pmatrix}
 1&k&\binom{k}{2}\\
 0&1&k\\
 0&0&1
 \end{pmatrix}.
\]
Liczby całkowite \(k\) i
\(\binom{k}{2}\) rozumiemy jako ich obrazy w \(K\);
w szczególności wzór obowiązuje również
w charakterystyce \(2\).
\end{przyklad}

\section{Funkcje analityczne macierzy zespolonych}
\label{sec:lin18-analytic-functions}

W tej sekcji pracujemy nad ciałem \(\mathbb C\).
Niech \(V\) będzie skończenie wymiarową
przestrzenią zespoloną, a
\(T\in\operatorname{End}_{\mathbb C}(V)\).
Dla \(\lambda\in\sigma_{\mathbb C}(T)\)
niech \(a_\lambda\) oznacza wykładnik
\((t-\lambda)\) w \(m_T\).
Przypomnijmy, że
\(N=T-\sum_\lambda\lambda P_\lambda\).

\begin{definicja}[Funkcja analityczna operatora]
\label{def:lin18-analytic-calculus}
Niech \(U\subset\mathbb C\)
będzie zbiorem otwartym zawierającym
\(\sigma_{\mathbb C}(T)\),
a \(f\colon U\to\mathbb C\)
funkcją holomorficzną.
Definiujemy
\[
 f(T):=
 \sum_{\lambda\in\sigma_{\mathbb C}(T)}
 \sum_{j=0}^{a_\lambda-1}
 \frac{f^{(j)}(\lambda)}{j!}\,N^jP_\lambda.
\]
Jest to skończona suma
endomorfizmów przestrzeni \(V\).
Dla \(V=0\) przyjmujemy \(f(T)=0\).
\end{definicja}

\begin{twierdzenie}[Własności rachunku funkcji analitycznych]
\label{thm:lin18-analytic-calculus}
Niech \(T\) będzie endomorfizmem
skończenie wymiarowej przestrzeni
zespolonej \(V\),
a \(f,g\) funkcjami holomorficznymi
na wspólnym otwartym otoczeniu
\(\sigma_{\mathbb C}(T)\).
Wówczas dla \(\alpha,\beta\in\mathbb C\)
\[
 \begin{gathered}
 (\alpha f+\beta g)(T)
    =\alpha f(T)+\beta g(T),\\
 (fg)(T)=f(T)g(T),\qquad
 1(T)=I_V,\qquad
 t(T)=T.
 \end{gathered}
\]
Jeżeli \(f(\lambda)\ne0\)
dla każdej \(\lambda\in\sigma_{\mathbb C}(T)\),
to \(f(T)\) jest odwracalny, a
\[
 f(T)^{-1}=(1/f)(T).
\]
Dla wielomianu \(p\in\mathbb C[t]\)
definicja \(p(T)\) zgadza się
ze zwykłym podstawieniem operatora
do wielomianu.
\end{twierdzenie}
\begin{proof}
Liniowość wynika bezpośrednio
z liniowości pochodnych.
Ustalmy \(\lambda\) i oznaczmy
\(W_\lambda=\operatorname{im}P_\lambda\)
oraz \(M=N|_{W_\lambda}\).
Mamy \(M^{a_\lambda}=0\).
W otoczeniu \(\lambda\)
szeregi Taylora \(f\) oraz \(g\)
spełniają
\[
 \begin{aligned}
 f(\lambda+u)
 &=\sum_{j=0}^{a_\lambda-1}
   \frac{f^{(j)}(\lambda)}{j!}u^j
   +u^{a_\lambda}r_f(u),\\
 g(\lambda+u)
 &=\sum_{j=0}^{a_\lambda-1}
   \frac{g^{(j)}(\lambda)}{j!}u^j
   +u^{a_\lambda}r_g(u),
 \end{aligned}
\]
gdzie \(r_f,r_g\) są holomorficzne
w pewnym otoczeniu zera.
Współczynniki przy potęgach \(u^j\)
dla \(j<a_\lambda\)
w rozwinięciu iloczynu \(fg\)
są dokładnie współczynnikami iloczynu
obciętych wielomianów po pominięciu
potęg \(u^{a_\lambda}\) i wyższych.
Podstawiając \(M\), otrzymujemy
\[
 (fg)(T)|_{W_\lambda}
 =f(T)|_{W_\lambda}\,g(T)|_{W_\lambda},
\]
ponieważ \(M^{a_\lambda}=0\).
Wszystkie rzuty przemieniają się z \(T\),
a różne składniki \(W_\lambda\)
tworzą sumę bezpośrednią.
Równość na każdym składniku
daje więc \((fg)(T)=f(T)g(T)\).

Funkcja stała \(1\) ma tylko
zerowy współczynnik Taylora poza
wyrazem stałym, zatem \(1(T)=\sum_\lambda
P_\lambda=I_V\).
Funkcja \(t\mapsto t\) ma
na \(W_\lambda\) postać
\(\lambda I+M=T|_{W_\lambda}\),
co daje \(t(T)=T\).
Dla wielomianów zgadza się to
z rachunkiem z
Twierdzenia~\ref{thm:lin18-polynomial},
gdyż ich współczynniki
w rozwinięciu przy \(\lambda\)
są pochodnymi dzielonymi przez silnie.

Jeżeli \(f(\lambda)\ne0\)
dla wszystkich wartości własnych,
to po ewentualnym zmniejszeniu
otoczeń poszczególnych punktów
funkcja \(1/f\) jest holomorficzna
na ich sumie; ta suma zawiera
całe widmo \(T\).
Stosując właśnie udowodnioną
multiplikatywność do \(f\) i \(1/f\),
otrzymujemy
\(f(T)(1/f)(T)=I_V=(1/f)(T)f(T)\).
\end{proof}

\begin{uwaga}[Funkcja nie musi być analityczna na całej płaszczyźnie]
\label{rem:lin18-domain}
Wystarcza, aby \(f\) była
holomorficzna w pewnym otwartym otoczeniu
skończonego zbioru wartości własnych \(T\).
Nie musi być określona na całej \(\mathbb C\).
Na przykład jeśli \(T\) jest odwracalny,
funkcja \(f(z)=1/z\) jest holomorficzna
w otoczeniu jego widma, a
Twierdzenie~\ref{thm:lin18-analytic-calculus}
daje \(f(T)=T^{-1}\).
\end{uwaga}

\begin{twierdzenie}[Wykładnicza operatora o postaci Jordana]
\label{thm:lin18-exponential}
Niech \(V\) będzie skończenie wymiarową
przestrzenią zespoloną i
\(T\in\operatorname{End}_{\mathbb C}(V)\).
Dla \(z\in\mathbb C\) zdefiniujmy
\[
 e^{zT}:=\sum_{k=0}^{\infty}
             \frac{z^kT^k}{k!}.
\]
Szereg zbiega w dowolnej normie operatorowej.
Przy oznaczeniach rozkładu
Jordana--Chevalleya zachodzi
\[
 e^{zT}=
 \sum_{\lambda\in\sigma_{\mathbb C}(T)}
 e^{z\lambda}
 \sum_{j=0}^{a_\lambda-1}
 \frac{z^j}{j!}\,N^jP_\lambda.
\]
Ponadto dla \(z,w\in\mathbb C\)
\[
 e^{0T}=I_V,\qquad
 e^{(z+w)T}=e^{zT}e^{wT},\qquad
 (e^{zT})^{-1}=e^{-zT}.
\]
\end{twierdzenie}
\begin{proof}
Wybierzmy dowolną normę na \(V\)
oraz indukowaną normę operatorową.
Dla każdego \(k\ge0\) mamy
\(\|T^k\|\le\|T\|^k\), zatem
\[
 \sum_{k=0}^{\infty}
 \left\|\frac{z^kT^k}{k!}\right\|
 \le
 \sum_{k=0}^{\infty}\frac{(|z|\|T\|)^k}{k!}
 =e^{|z|\|T\|}<\infty.
\]
Stąd zbieżność bezwzględna szeregu
w przestrzeni operatorów.
W przypadku \(V=0\) mamy jedyną
zerową normę operatorową i teza
jest zgodna z konwencjami.

Na \(W_\lambda=\operatorname{im}P_\lambda\)
operator \(T\) jest równy \(\lambda I+N\)
z \(N^{a_\lambda}=0\).
Ponieważ \(\lambda I\) i \(N\) przemieniają się,
z dwumianu Newtona w skończonych sumach
oraz z bezwzględnej zbieżności szeregu
wynika
\[
 \begin{aligned}
 e^{zT}|_{W_\lambda}
 &=e^{z\lambda}\,
   \sum_{j=0}^{a_\lambda-1}
          \frac{z^j}{j!}N^j|_{W_\lambda}.
 \end{aligned}
\]
Jest to też bezpośrednio definicja
funkcji \(f(\zeta)=e^{z\zeta}\)
z Definicji~\ref{def:lin18-analytic-calculus}.
Sumując po przestrzeniach \(W_\lambda\),
otrzymujemy żądany wzór.

Tożsamość \(e^{0T}=I_V\)
wynika z definicji.
Z
Twierdzenia~\ref{thm:lin18-analytic-calculus}
i równości skalarnych
\(e^{(z+w)\zeta}=e^{z\zeta}e^{w\zeta}\)
otrzymujemy
\(e^{(z+w)T}=e^{zT}e^{wT}\).
Podstawiając \(w=-z\), uzyskujemy
\((e^{zT})^{-1}=e^{-zT}\).
\end{proof}

\begin{przyklad}[Rozwiązanie liniowego układu różniczkowego]
\label{prz:lin18-ode}
Niech \(T=J_3(2)\) działa na \(\mathbb C^3\).
Wówczas
\(T=2I+N\), \(N^3=0\), a
\[
 e^{tT}
 =e^{2t}
 \begin{pmatrix}
 1&t&t^2/2\\
 0&1&t\\
 0&0&1
 \end{pmatrix}
 \qquad(t\in\mathbb R).
\]
Rozważmy zagadnienie początkowe
\(x'(t)=Tx(t)\), \(x(0)=e_3\).
Funkcja
\[
 x(t)=e^{tT}e_3
     =e^{2t}\begin{pmatrix}t^2/2\\t\\1\end{pmatrix}
\]
spełnia warunek początkowy.
Różniczkując współrzędne, otrzymujemy
\[
 x'(t)
 =e^{2t}
 \begin{pmatrix}t^2+t\\2t+1\\2\end{pmatrix}.
\]
Z drugiej strony
\[
 Tx(t)
 =e^{2t}
 \begin{pmatrix}t^2+t\\2t+1\\2\end{pmatrix},
\]
więc równanie rzeczywiście zachodzi.
Wzór pokazuje, dlaczego przy blokach Jordana
oprócz czynników wykładniczych
pojawiają się czynniki wielomianowe.
\end{przyklad}

\section{Podsumowanie}
\label{sec:lin18-summary}

Dla operatora o rozszczepionym wielomianie
minimalnym otrzymaliśmy wielomianowe rzuty
\(P_\lambda\) na uogólnione przestrzenie własne.
Są one kanoniczne, przemieniają się z operatorem,
a ich suma jest identycznością, choć nie muszą
być ortogonalne.

Rozkład Jordana--Chevalleya
\(T=S+N\), \(SN=NS\), oddziela
część diagonalizowalną
\(S=\sum_\lambda\lambda P_\lambda\)
od nilpotentnej \(N=T-S\).
Obie są wielomianami w \(T\),
a rozkład jest jednoznaczny.
Wielomiany w \(T\), jego potęgi
i odwrotność można obliczać osobno
na każdej uogólnionej przestrzeni własnej,
używając skończonych rozwinięć wokół
wartości własnych.

Nad \(\mathbb C\) ten sam mechanizm
określa funkcje analityczne macierzy.
W szczególności wykładnicza macierzy
zachowuje wzór na wykładniczą części
diagonalizowalnej, uzupełniony
o skończone wielomiany części nilpotentnej.
Zastosowaliśmy go do rozwiązania
liniowego układu różniczkowego.
